% !TEX root = ../../tree_notes.tex
% =====================================================================
%  Chapter 2 opener -- "Chapter at a Glance", reading paths + ladder
%  staircase, running example, roadmap, topic table, recognition guide,
%  notation.
%  NO numbered boxes here (they would be numbered 2.0.x).
% =====================================================================

\begin{chaptertheorybox}{Chapter at a Glance}
A \term{binary tree} is a rooted tree in which every node has at most two children,
distinguished as \emph{left} and \emph{right}. It is the shape interviews and the
lighter contest problems are built from, and the data structure every later chapter
either balances, indexes, or decomposes. But this chapter is not a catalogue of
binary-tree trivia -- it is about the \emph{one idea} from which all of them flow: a
binary tree is a node plus two smaller binary trees, so
\textbf{every question about a tree is one function call on the node plus two
function calls on the children}. We organise the whole world of tree recursion by a
single question: \emph{which way does the information travel?} Values carried
\emph{down} the calls give the \term{top-down} pattern; values returned \emph{up}
give the \term{bottom-up} pattern; and the hardest problems combine both, or rebuild
the tree itself while recursing.

\textbf{You will learn to}
\begin{itemize}[nosep]
  \item read and write the \code{TreeNode} pointer structure and recurse over it
        without double-visiting or forgetting the \nil{} base case;      \hfill\lv{1}
  \item decide, from the problem's phrasing, whether state travels \emph{down}
        (path, depth, ancestor context) or answers travel \emph{up}
        (subtree sums, heights, aggregates);                            \hfill\lv{2}
  \item write the top-down skeleton with a carried parameter, and recognise when it
        degenerates into $O(n^2)$ and needs prefix state;               \hfill\lv{3}--\lv{4}
  \item write the bottom-up skeleton, return \emph{two} values from one recursion, and
        keep a global answer updated on the way up;                     \hfill\lv{3}--\lv{5}
  \item transform tree structure in place -- rewire child pointers, prune, delete and
        split -- with recursion that returns the new root of each subtree. \hfill\lv{3}--\lv{4}
\end{itemize}

\textbf{Prerequisites.} Chapter~1: the vocabulary of roots, leaves, subtrees, depth
and height (\secref{c01-anatomy}), and the habit of structural recursion shown on
Same Tree (\secref{c01-what-is-tree}). Comfort with $\cpp$ pointers.

\textbf{In contests and interviews.} Binary-tree recursion is the single most tested
interview family. Once you own the two directions -- down and up -- roughly every
$\lv1$--$\lv4$ LC tree problem becomes the same twelve-line function with a
different combine step. This chapter supplies that frame once; Chapters
\chref{c03}--\chref{c05} and the DP chapters reuse it without re-teaching it.
\end{chaptertheorybox}

% ---------------- Station 1b: reading paths + ladder position ----------------
\begin{chaptertheorybox}{How to Read This Chapter (two paths up the ladder)}
Level badges $\lv1$--$\lv6$ work as in Chapter~1: $\lv1$ is \emph{``what does this
function return''}, $\lv3$ is the solid interview medium, $\lv5$ is LC Hard with a
trick on top.

\textbf{Beginner path ($\lv1$--$\lv3$).} Read \secref{c02-model} completely -- it
fixes the node picture, the base-case discipline and the running example of the
chapter. Then take \secref{c02-topdown} and \secref{c02-bottomup}: definitions,
properties, skeletons and every $\lv1$--$\lv3$ solved problem. You may \emph{skip on
the first pass}: the prefix-hashmap rung of Path Sum III, the two-channel
House Robber III build, the $\lv4$+ variation entries, the hardening notes and the
chapter's Summit. Do every recall card. At the end of this path you can solve every
$\lv3$ binary-tree problem in this chapter.

\textbf{Full path.} Everything, in order -- it ends with a scaled summit (LC Hard
968 Binary Tree Cameras: the hardest problem solvable with this chapter's bottom-up
tools alone) and the reduction gallery showing how every lab problem is the same
skeleton relabelled. The true $\lv5$--$\lv6$ summits of the recursion idea arrive
with tree DP (\chref{c10}) and rerooting (\chref{c20}).
\end{chaptertheorybox}

\begin{figure}[H]
\centering
\begin{tikzpicture}[diagram, x=9mm, y=7mm]
  \foreach \st/\lab in {1/lv1,2/lv2,3/lv3,4/lv4,5/lv5,6/lv6} {
    \fill[\lab!22, draw=\lab, line width=0.8pt] (\st-1,0) rectangle (\st,\st*0.52);
    \node[font=\scriptsize\bfseries, text=\lab] at (\st-0.5,\st*0.52-0.42) {L\st};}
  \node[annot, anchor=south west, font=\scriptsize]
    at (0.05,0.62) {2.1 the node};
  \node[annot, anchor=south west, font=\scriptsize]
    at (1.05,1.14) {2.2 top-down};
  \node[annot, anchor=south west, font=\scriptsize]
    at (2.05,1.66) {2.3 bottom-up};
  \node[annot, anchor=south west, font=\scriptsize]
    at (3.05,2.18) {2.4 transforms};
  \node[annot, anchor=south west, font=\scriptsize, text=lv5]
    at (4.05,2.70) {scaled summit (968)};
  \node[annot, anchor=south west, font=\scriptsize, text=lv6]
    at (5.05,3.22) {true summits (tree DP, ch.\,10)};
\end{tikzpicture}
\caption*{Chapter~2 on the six-step ladder: the four topics carry you to $\lv4$, the
scaled summit (Cameras) sits on step~5, and the genuine $\lv6$ recursion problems
wait in the tree-DP chapters. Beginner path $=$ steps 1--3.}
\end{figure}

% ---------------- The running example ----------------
\subsection*{The Running Example $R$}

Every figure, worked example and trace in this chapter uses the \emph{same} six-node
tree $R$. Learn its shape once and you can follow every trace without recomputing
new data.

\begin{figure}[H]
\centering
\begin{tikzpicture}[diagram]
  \node[troot] (r) at (0,0) {$1$};
  \node[tnode] (a) at (-1.7,-1.4) {$2$};
  \node[tnode] (b) at (1.7,-1.4) {$3$};
  \node[tleaf] (c) at (-2.6,-2.8) {$4$};
  \node[tleaf] (d) at (-0.8,-2.8) {$5$};
  \node[tnode] (e) at (1.7,-2.8) {$6$};
  \draw[tedge] (r) -- (a); \draw[tedgehl] (r) -- (b);
  \draw[tedge] (a) -- (c); \draw[tedge] (a) -- (d);
  \draw[tedgehl] (b) -- (e);
  \node[annot, anchor=west] at (-3.4,-2.1) {node $2$:};
  \node[annot, anchor=west, align=left, font=\scriptsize] at (-3.4,-2.55) {$\depth=1$, both children,\\not a leaf};
  \node[annot, anchor=east, align=left, font=\scriptsize] at (3.4,-1.4) {node $3$:\\$\depth=1$, $\height=1$,\\right child only};
\end{tikzpicture}
\caption{The running example $R$. Values double as labels. Root $1$; left child $2$
with children $4,5$; right child $3$ with right child $6$ (a left child \emph{missing}
on purpose -- one-child nodes are where base-case bugs surface). Leaves $\{4,5,6\}$;
$\height(R)=2$ edges; the highlighted chain $1\to3\to6$ is a root-to-leaf path, as
are $1\to2\to4$ and $1\to2\to5$.}\label{fig:c02-running}
\end{figure}

Facts about $R$ you will re-use in the traces (and should be able to re-derive
instantly): the three root-to-leaf paths are $1\text{-}2\text{-}4$,
$1\text{-}2\text{-}5$, $1\text{-}3\text{-}6$; read as decimal numbers they are
$124+125+136=385$; the node values along the leftmost path sum to $7$; and the
subtree at $2$ has $\size=3$ with sum $11$, while the subtree at $3$ has $\size=2$
with sum $9$.

% ---------------- Roadmap diagram ----------------
\begin{figure}[H]
\centering
\begin{tikzpicture}[diagram, node distance=4mm,
    top/.style={boxnode, fill=primaryblue!10, draw=primaryblue, text width=25mm,
                minimum height=9mm, font=\small},
    res/.style={annot, text width=25mm, anchor=north}]
  \node[top] (t1) {\textbf{2.1} The node};
  \node[top, right=of t1] (t2) {\textbf{2.2} Top-down};
  \node[top, right=of t2] (t3) {\textbf{2.3} Bottom-up};
  \node[top, right=of t3] (t4) {\textbf{2.4} Transforms};
  \foreach \a/\b in {t1/t2, t2/t3, t3/t4} \draw[flow] (\a) -- (\b);
  \node[res] at ([yshift=-2mm]t1.south) {\code{TreeNode}, base cases,
    structural recursion on $R$};
  \node[res] at ([yshift=-2mm]t2.south) {carry state down: paths, ancestors,
    running sums, prefixes};
  \node[res] at ([yshift=-2mm]t3.south) {return values up: subtree aggregates,
    heights, two-channel recursion};
  \node[res] at ([yshift=-2mm]t4.south) {rewire pointers: invert-family, pruning,
    delete-and-split};
\end{tikzpicture}
\caption*{The four topics, left to right, each with the result you take away.
2.1 builds the machine; 2.2 and 2.3 are its two directions of data flow; 2.4 points
the machine at the tree itself.}
\end{figure}

% ---------------- Topic table ----------------
\begin{center}
\renewcommand{\arraystretch}{1.25}
\begin{tabular}{@{}p{0.26\linewidth} p{0.36\linewidth} c p{0.22\linewidth}@{}}
\toprule
\textbf{Topic} & \textbf{Key results} & \textbf{Levels} & \textbf{Solved problems} \\
\midrule
\secref{c02-model}\ \ The node and structural recursion
  & \code{TreeNode}; base-case discipline; visit-count $=$ $O(n)$ audit
  & $\lv1$--$\lv2$
  & \probref{c02-root-sum}, \probref{c02-second-min}, \probref{c02-paths} \\
\addlinespace
\secref{c02-topdown}\ \ Top-down: carrying state down
  & path-list accumulation; prefix-count trick on trees
  & $\lv2$--$\lv4$
  & \probref{c02-path-sum-ii}, \probref{c02-sum-numbers}, \probref{c02-grandparent},
    \probref{c02-anc-diff}, \probref{c02-path-sum-iii} \\
\addlinespace
\secref{c02-bottomup}\ \ Bottom-up: returning values up
  & aggregate-and-answer split; two-channel return
  & $\lv2$--$\lv5$
  & \probref{c02-tilt}, \probref{c02-prune}, \probref{c02-rob},
    \probref{c02-coins}, \probref{c02-deepest}, \probref{c02-zigzag} \\
\addlinespace
\secref{c02-transform}\ \ Structural transformation
  & ``return the new root of this subtree'' discipline
  & $\lv3$--$\lv4$
  & \probref{c02-print}, \probref{c02-create-tree}, \probref{c02-flip-eq},
    \probref{c02-delete-forest} \\
\bottomrule
\end{tabular}
\end{center}

% ---------------- Problem recognition guide ----------------
\subsection*{Problem Recognition Guide --- When You See \emph{this}, Think \emph{that}}

Before you code, ask \emph{which way the information moves}. Match the problem's
phrasing against this table; every row is a pattern formally built in the section
named at the right.

\begin{center}
\renewcommand{\arraystretch}{1.3}
\begin{tabular}{@{}p{0.40\linewidth} p{0.38\linewidth} l@{}}
\toprule
\textbf{When the problem says\ldots} & \textbf{\textcolor{deepblue}{\textbf{Think\ldots}}} & \textbf{Section} \\
\midrule
``list every root-to-leaf path / number''; ``for each node, using its ancestors''
  & carry a parameter \emph{down}: the path so far, the ancestor context & \secref{c02-topdown} \\
\addlinespace
``paths \emph{through} a node with a sum property, node not on the path''
  & top-down is not enough -- carry \emph{prefix} state and count matches & \secref{c02-topdown} \\
\addlinespace
``sum / count / max over each subtree''; ``tilt, prune, distribute''
  & \emph{return} the aggregate \emph{up}; update the answer at each node & \secref{c02-bottomup} \\
\addlinespace
``the answer at $v$ needs the subtree sum \emph{and} the best answer below''
  & two-channel recursion: return a pair, decide at the parent & \secref{c02-bottomup} \\
\addlinespace
``invert, prune, delete nodes, rebuild, flip to match''
  & recursion that \emph{returns the new root} of each subtree & \secref{c02-transform} \\
\addlinespace
``is this node the root such that\ldots'' / global optimum over subtrees
  & bottom-up with a candidate check at every node, not a fresh DFS per node & \secref{c02-bottomup} \\
\bottomrule
\end{tabular}
\end{center}

\begin{center}
\renewcommand{\arraystretch}{1.25}
\begin{tabular}{@{}p{0.26\linewidth} p{0.46\linewidth} p{0.20\linewidth}@{}}
\toprule
\textbf{Constraint you see} & \textbf{What it permits} & \textbf{Chapter-2 tool} \\
\midrule
$n \le 10^{4}$ (typical LC)  & even two full DFS passes per node fit -- do not over-optimise & \secref{c02-model} \\
$n \le 10^{5}$               & exactly one visit per node, $O(1)$ combine -- the skeletons of 2.2/2.3 & \secref{c02-bottomup} \\
values up to $10^{4}$, sums over paths & the running sum fits \code{int}, but the \emph{count} of paths may not -- check before using \code{long long} & \secref{c02-topdown} \\
tree may be a chain of $10^{5}$ nodes & recursion depth is the danger -- iterative form or raised stack & \secref{c02-model} \\
\bottomrule
\end{tabular}
\end{center}

\begin{notebox}[Not this chapter?]
If the children are \emph{ordered by value} and the tree must stay that way for
lookup, you want BSTs (\chref{c05}). If the problem names an order -- ``inorder'',
``level order'', ``zigzag level'' -- that is traversal machinery (\chref{c03}). If
the tree is stored in an array because it is a heap, go to \chref{c04} and
\chref{c07}. This chapter owns the recursion, not the bookkeeping of visits.
\end{notebox}

% ---------------- Notation ----------------
\subsection*{Notation Used in This Chapter}

\begin{center}
\renewcommand{\arraystretch}{1.2}
\begin{tabular}{@{}p{0.20\linewidth} p{0.28\linewidth}@{\quad}p{0.16\linewidth} p{0.28\linewidth}@{}}
\toprule
\textbf{Symbol} & \textbf{Meaning} & \textbf{Symbol} & \textbf{Meaning} \\
\midrule
$T$, $R$   & a binary tree; $R$ = the running example & $\depth(v)$  & edges from root down to $v$ \\
$v\!\to\!\mathit{left}$, $v\!\to\!\mathit{right}$
           & the two child pointers of $v$ & $\height(v)$  & edges from $v$ to deepest leaf below \\
$\nil$     & the empty / \code{null} child pointer & $\size(v)$    & nodes in the subtree $\subtree(v)$ \\
$\subtree(v)$ & the tree rooted at $v$ & $\Sigma(v)$    & sum of values in $\subtree(v)$ \\
$n$        & number of nodes & $h$            & height of the whole tree \\
$\mathit{ans}$  & global answer updated during a DFS & $S$  & the carried path-sum (2.2) \\
$\mathit{cnt}[s]$ & hashmap of root-prefix sums $s$ & $\delta(v)$ & parity flag (2.2) \\
\bottomrule
\end{tabular}
\end{center}

As in Chapter~1, $\height$ and $\depth$ are counted in \textbf{edges}
($\height(\nil)=-1$), and every listing is complete $\cpp$ with the standard
\code{TreeNode} struct, so it compiles against the LeetCode driver as-is.
